\documentclass[12pt]{article}
\usepackage[utf8]{inputenc}
\usepackage[T1]{fontenc}
\usepackage[spanish]{babel}
\usepackage{geometry}
\usepackage{amsmath,amssymb}
\usepackage{enumitem}

\geometry{margin=2.5cm}

\begin{document}

\begin{center}
{\Large Transcripción cruda}\\[2mm]
{\large 28May26}
\end{center}

\section*{Fotos 1--6}

\textit{Lema.} Sea $(X,d)$ un espacio métrico.
Sean $a\in X$ y $r>0$.
Definimos
\[
V_r(a):=\{x\in X\mid d(x,a)\le r\}.
\]
Ojo:
\[
N_r(a)=\{x\in X\mid d(x,a)<r\}\subset V_r(a).
\]
Entonces:
\begin{enumerate}[label=\roman*)]
    \item $V_r(a)$ es cerrado.
    \item $N_r(a)\subset V_r(a)$, aunque la contención puede ser estricta.
    \item Si $X=\mathbb R^k$, entonces $\overline{N_r(a)}=V_r(a)$.
\end{enumerate}

\textit{Dem.} i) Demostraremos que $(V_r(a))^c$ es abierto.
Sea $b\in (V_r(a))^c$.
Sea
\[
d_0=\frac12\bigl(d(b,a)-r\bigr)>0.
\]
Afirmamos que
\[
N_{d_0}(b)\subset (V_r(a))^c.
\]
Si $x\in N_{d_0}(b)$, entonces
\[
r<d(a,b)\le d(a,x)+d(x,b)\le d(a,x)+d_0,
\]
de donde $d(a,x)>r$.
Por tanto $x\notin V_r(a)$.

ii) es inmediato.

Para mostrar que la contención puede ser estricta, consideremos la métrica discreta
\[
d(x,y)=\begin{cases}
0,&x=y,\\
1,&x\ne y.
\end{cases}
\]
Entonces
\[
N_1(a)=\{x\in X\mid d(a,x)<1\}=\{a\},
\qquad
V_1(a)=\{x\in X\mid d(a,x)\le 1\}=X.
\]

iii) Si $X=\mathbb R^k$, entonces $\overline{N_r(a)}=V_r(a)$.
Sea $x\in V_r(a)\setminus \overline{N_r(a)}$.
Tomemos
\[
x_\lambda=(1-\lambda)a+\lambda x,\qquad 0\le \lambda\le 1.
\]
Entonces
\[
d(a,x_\lambda)=\|a-((1-\lambda)a+\lambda x)\|
=\|\lambda(a-x)\|
=|\lambda|\,\|a-x\|
\le \lambda r.
\]
Si $0\le\lambda<1$, se tiene $d(a,x_\lambda)<r$, así que $x_\lambda\in N_r(a)$.
Al dejar $\lambda\to 1$, se obtiene $x\in\overline{N_r(a)}$.

\section*{Fotos 7--11}

\textit{Ejemplo.} En la métrica discreta,
\[
N_1(a)=\{a\}\qquad\text{y}\qquad V_1(a)=X.
\]
Así que la inclusión $N_1(a)\subset V_1(a)$ puede ser estricta.

\textit{Observación.} En $\mathbb R^k$ la vecindad cerrada coincide con la cerradura de la abierta.
Esto se ve porque si $x\in V_r(a)$, entonces el segmento
\[
(1-\lambda)a+\lambda x,\qquad 0\le \lambda<1,
\]
permanece dentro de $N_r(a)$ y se aproxima a $x$ cuando $\lambda\to 1$.

\section*{Fotos 12--13}

\textit{Teo.} Sea $T\subset\mathbb R^d$ infinito y acotado.
Entonces $T$ cabe en una $d$-celda,
\[
T\subset [a_1,b_1]\times\cdots\times[a_d,b_d].
\]
Ya sabemos que
\[
[a_1,b_1]\times\cdots\times[a_d,b_d]
\]
es compacto.
Por tanto, $T$ es un subconjunto infinito de un compacto; luego tiene punto límite en $\mathbb R^d$.

\textit{Corolario (Weierstrass).} Todo subconjunto infinito acotado de $\mathbb R^d$ tiene un punto límite en $\mathbb R^d$.

\textit{Normas en $\mathbb R^k$.} Para $x,a\in\mathbb R^k$:
\[
\|x-a\|_1=\sqrt{\sum_{i=1}^k(x_i-a_i)^2},
\qquad
\|x-a\|_2=\max\{|x_i-a_i|:i=1,\dots,k\}.
\]
Además,
\[
\|x\|_2=\|x-0\|_2=\max\{|x_1|,\dots,|x_k|\},
\qquad
\|x\|_1=\sqrt{\sum_{i=1}^k x_i^2}.
\]

\textit{Ejercicio.} Probar lo siguiente:
\begin{enumerate}[label=\roman*)]
    \item Si $x\in\mathbb R^k$ y $\alpha\in\mathbb R$, entonces
    \[
    \|\alpha x\|_2=|\alpha|\,\|x\|_2.
    \]
    \item
    \[
    \|x\|_2=0 \iff x=0.
    \]
    \item
    \[
    \|x+y\|_2\le \|x\|_2+\|y\|_2.
    \]
\end{enumerate}

\end{document}
